\subsection{I-36. A pairwise dominant center preserves maximal minors; five pairs for singular width-three words}
\label{I-36}

\paragraph{Statement and scope.}
Fix distinct Lagrange nodes $a_1,\ldots,a_n$ in a working field
$E$, and let $\lambda_i(Y)$ be their degree-at-most-$(n-1)$
Lagrange polynomials. The coefficient field is a subfield of $E$.
For $n\ge1$, define
\[
 \alpha_i=t+z\lambda_i(y),\qquad
 \mathcal G_i^{(5)}=t+\sum_{j=1}^{5}z_j\lambda_i(y_j).
 \tag{I.84}
\]
All displayed parameters are algebraically independent.

The map $\boldsymbol\alpha$ preserves every nonzero maximal-size
minor of every interval in a univariate matrix word whose individual
factors have generic rank at most that size. More explicitly, if
for integers $1\le r\le w$, the $A_i(X_i)$ are $w\times w$
polynomial matrices of generic rank at
most $r$, with distinct variables at the different positions, then
every nonzero $r\times r$ minor of every interval product remains
nonzero after $\mathbf X=\boldsymbol\alpha$. Identically zero minors
remain zero. The map depends on neither the matrices nor their order.

Consequently $\mathcal G^{(5)}$ hits every nonzero width-at-most-three
ROABP admitting a representation in which every variable-dependent
matrix factor has generic rank at most two. Constant invertible
width-three factors may be absorbed into neighboring factors or
the boundaries. Arbitrarily many singular factors and arbitrary
drops in the ranks of longer intervals are allowed. The variable
order is unknown to the sampler. This improves I-30's six-pair bound
on this specified singular-factor subclass; it does not prove that
two centered pairs suffice, and it does not extend this improvement
to arbitrary variable-dependent invertible width-three factors.

\paragraph{1. Dominance on every coordinate pair.}
For $i\ne j$, the polynomial $\lambda_i(y)-\lambda_j(y)$ is nonzero:
its value at $y=a_i$ is one. Over $E(y)$, the linear substitution
\[
 \begin{pmatrix}X_i\\X_j\end{pmatrix}
 =\begin{pmatrix}1&\lambda_i(y)\\1&\lambda_j(y)\end{pmatrix}
   \begin{pmatrix}t\\z\end{pmatrix}
 \tag{I.85}
\]
is invertible. It therefore induces an injection
$E[X_i,X_j]\hookrightarrow E(y)[t,z]$. Every nonzero polynomial
depending on at most two specified coordinates stays nonzero.
The singleton case follows directly from the coefficient one on
$t$; it covers $n=1$. The same argument holds after any extension
of the coefficient field independent of the new parameters.

This is a formal polynomial assertion. It does not require a
single numerical parameter choice to work over a small base field;
the later interpolation step supplies the needed finite evaluations.

\paragraph{2. Rank-one words factor into adjacent bivariate contractions.}
Let $B_i(X_i)$ be nonzero polynomial matrices of generic rank one.
Choose any nonzero entry $b_i$ of $B_i$, let $p_i$ be its column,
and let $q_i$ be its row. The vanishing of the two-by-two minors gives
\[
 B_i=\frac{p_iq_i}{b_i}
 \quad\hbox{over the rational function field.}
\]
For every interval $[s,t]$ and constant boundary row $u$ and column
$v$ of the appropriate sizes, clearing the nonzero univariate
pivots gives the exact polynomial identity
\[
 \left(\prod_{i=s}^{t}b_i\right)
       uB_sB_{s+1}\cdots B_tv
 = (up_s)(q_tv)
       \prod_{i=s}^{t-1}(q_ip_{i+1}).
 \tag{I.86}
\]
The two endpoint factors are univariate. Each middle factor depends
on only the two distinct adjacent variables $X_i,X_{i+1}$.
This remains valid for a one-factor interval, with an empty middle
product. The variable sets of different factors on the right may
overlap; their nonzero product is still nonzero in a polynomial
domain.

If the original scalar contraction is nonzero, every factor on the
right of (I.86) is nonzero. Part~1 preserves all those factors and
all the pivots under the same substitution $\boldsymbol\alpha$.
The substituted scalar contraction is therefore nonzero. In
particular every nonzero entry of every rank-one interval word is
preserved. This proof uses scalar product closure after an exact
rank-one factorization, not an unsupported composition of arbitrary
constant-boundary tests with shared parameters.

\paragraph{3. Applying the factorization to maximal minors.}
For a matrix $A$, write $C_r(A)$ for its $r$th compound matrix.
Cauchy--Binet gives
\[
 C_r(A_s\cdots A_t)=C_r(A_s)\cdots C_r(A_t).
 \tag{I.87}
\]
If a factor has generic rank below $r$, its compound is zero and
every interval containing it has all $r$-minors identically zero.
Otherwise all factors in the interval have generic rank exactly
$r$, so each $C_r(A_i)$ has generic rank one. Part~2 applies to
this compound word and preserves each nonzero entry of (I.87),
which is precisely the corresponding $r$-minor of the original
interval. This proves the first assertion of the statement.

Compound matrices occur only in the proof. The sampler does not
construct them, recover their pivot entries, inspect the input
representation, or enumerate the unknown variable order. All
intervals and minors are preserved by the same formal map: the
argument establishes a nonzero image for each one individually,
and any finite product of these nonzero images is nonzero.

\paragraph{4. Completing the rank-one stage in width three.}
Remove constant invertible factors from the width-three word by
absorbing them into neighboring matrices or the constant boundary
vectors. This does not increase the generic rank of any remaining
factor. Unused variables cause no problem for (I.85). A representation
of smaller width may be padded with zero rows and columns to width
three, so its factors also have generic rank at most two.

Apply parts~1--3 with $r=2$. Every nonzero two-by-two interval
minor survives at the center $\boldsymbol\alpha$. There are no
generic rank-three intervals in the remaining word. The two
constant rank-one sentinels used in I-21 to encode the scalar
output introduce no rank-three interval or additional variable.

Now use the local rank-test construction proved in I-21 and
the endpoint absorption of I-30. In each selected nonzero
rank-one test, the interior intervals have generic ranks at
least two. Here their ranks are exactly two, and they have
already been preserved at $\boldsymbol\alpha$. Every nonzero
univariate polynomial endpoint vector is nonzero at its center:
each $\alpha_i$ is transcendental over the original coefficient
field, since its coefficient on $t$ is one. Constant nonzero
endpoints are covered directly.

Pass to the field of the first-stage parameters and translate
to independent residual variables. I-30 absorbs the endpoints
into unit factors, producing a word with matching generic and
constant interval ranks in $\{2,3\}$. The complete support
lemma I-29 with $(w,r)=(3,2)$ supplies a nonzero coefficient
whose monomial has support at most
\[
 2(3-2)+3-1=4,
 \tag{I.88}
\]
including both endpoint positions. Four fresh SV pairs preserve
every selected rank-one test, because a polynomial with such a
nonzero low-support coefficient cannot vanish under $\mathrm{SV}_4$.
The new pairs are independent of the first-stage parameter field.
The earlier two-by-two minors remain nonzero, as is seen by setting
the new amplitudes to zero. The local-rank induction of I-21 now
preserves the nonzero scalar output. Adding the first pair and
the four residual pairs gives exactly $\mathcal G^{(5)}$ in (I.84).

This last argument is the already checked rank-one completion
from I-30, with its rank-two-preserving center supplied by the
new rank-one-compound factorization. It does not assert that
the first-stage center alone preserves rank-one intervals.

\paragraph{5. Exact finite sampling and field requirements.}
For a total-degree-at-most-$D$ input, the substituted polynomial
has degree at most $D$ in $t$, at most $(n-1)D$ in each of the
five selectors, and total degree at most $D$ in the five
amplitudes. The tensor grid in $t$ and the selectors, together
with the Newton simplex in the amplitudes used in I-21, yields
at most
\[
 (D+1)\bigl((n-1)D+1\bigr)^5\binom{D+5}{5}
 \tag{I.89}
\]
evaluations. The working field must contain at least
$\max\{n,D+1,(n-1)D+1\}$ distinct elements. If it is a proper
extension of the coefficient field, black-box evaluation over
that extension is part of the evaluation model. All arguments,
including the simplex interpolation, hold in every characteristic.
The constant-polynomial case $D=0$ needs only one evaluation;
the displayed formula also equals one. No polynomial dependence
on arbitrary width is inferred from this fixed-width result.
\hfill$\square$

\paragraph{Provenance and the remaining gate.}
The rank-one pivot factorization and scalar product closure are
the established mechanism recorded in I-B06 and P-14. The new
deduction here is that pairwise dominance supplies an entire
maximal-minor preservation stage with one centered interpolation
pair, followed by the previously proved four-pair completion.
No historical priority claim is made. The centered two-pair
question for rank-two reset words remains open; the explicit
uncentered failure P-16 does not decide it.
